\documentclass[11pt]{amsart}
\usepackage[T1]{fontenc}
\usepackage{lmodern,microtype,amsmath,amssymb,amsthm,mathtools}
\usepackage[margin=1.1in]{geometry}
\usepackage[hidelinks]{hyperref}
\usepackage{xurl}

\newtheorem{theorem}{Theorem}[section]
\newtheorem{lemma}[theorem]{Lemma}
\newtheorem{proposition}[theorem]{Proposition}
\theoremstyle{remark}
\newtheorem{remark}[theorem]{Remark}
\numberwithin{equation}{section}

\newcommand{\E}{\mathbb{E}}
\newcommand{\PP}{\mathbb{P}}

\hypersetup{pdftitle={Complete mixability of unimodal cubic densities at the boundary mean},pdfauthor={Yitzchak Shmalo}}

\begin{document}

\title[Unimodal cubic densities at the boundary mean]{Complete mixability of unimodal cubic densities at the boundary mean}
\author{Yitzchak Shmalo}
\address{Einstein Institute of Mathematics, The Hebrew University of Jerusalem, Jerusalem, Israel}
\email{yitzchak.shmalo@gmail.com}
\date{}
\subjclass[2020]{60E05, 60E15, 91G70}
\keywords{complete mixability, unimodal density, polynomial density, mean condition, dependence}

\begin{abstract}
A distribution is $3$-completely mixable if it is the common distribution
of three random variables whose sum is constant. We show that every unimodal
probability density on $[0,1]$ that is a polynomial of degree at most three
and has mean $1/3$, the smallest mean the mean condition allows, is
$3$-completely mixable. Such a density is either decreasing, and then
completely mixable by a theorem of Wang and Wang, or a convex combination of
two decreasing densities and one explicit $3$-completely mixable density with
an interior mode. Every $3$-completely mixable density $f$ on $[0,1]$ with mean
$1/3$ that is continuous at $0$ and $1$ satisfies $f(0)\ge4f(1)$; the norm
condition of Wang gives $f(0)\ge2f(1)$. A unimodal quartic density, strictly
positive on $[0,1]$, violates the weaker inequality, so the degree bound
cannot be raised. A quadratic density violates only the stronger one, so
unimodality cannot be dropped.
\end{abstract}

\maketitle

\section{Introduction}

A distribution $F$ on $\mathbb{R}$ is \emph{$n$-completely mixable} ($n$-CM)
if there are random variables $X_1,\dots,X_n$, each with distribution $F$,
such that $X_1+\dots+X_n$ is almost surely constant. We identify a
distribution with its distribution function, and we call a density $n$-CM
when the distribution it defines is $n$-CM. The notion was introduced by
Wang and Wang~\cite{WW11}; uniform distributions had already been shown to
be $n$-CM for $n\ge2$ by Gaffke and R\"uschendorf~\cite{GR81}
(see~\cite[Section~1.4]{Wan15}). Complete mixability enters the bounds on
the distribution of a sum of dependent risks with fixed marginals; see the
survey of Wang~\cite{Wan15} for the background and applications.

If $F$ is $n$-CM and has mean $\mu$, the constant sum is $n\mu$. If $F$ is
supported on $[a,b]$, with $a$ and $b$ the ends of its essential support,
then the \emph{mean condition}
\begin{equation}\label{eq:meancond}
a+\frac{b-a}{n}\le\mu\le b-\frac{b-a}{n}
\end{equation}
is necessary for $F$ to be $n$-CM~\cite[Proposition~2.1(7)]{WW11}. Wang and
Wang~\cite[Theorem~2.4 and Corollary~2.9]{WW11} proved that it is also
sufficient when $F$ has a monotone density on $[a,b]$. Symmetric unimodal
densities were studied for $n=3$ by Knott and Smith~\cite{KS06}, and
R\"uschendorf and Uckelmann~\cite{RU02} showed that a symmetric unimodal
density is $n$-CM for $n\ge2$; see~\cite[Theorem~2.2]{WW11} for the statement
in this form.
Unimodality cannot be dropped there: Yin and Zhu~\cite[Example~2.3]{YZ17}
show that the symmetric density proportional to $x^{2r}$ on $[-1,1]$,
$r\ge1$, is not $3$-CM.
Puccetti, Wang and Wang proved that a concave density on a bounded interval
is $n$-CM for $n\ge3$~\cite[Theorem~4.2]{PWW12}, and that a density bounded
below by $3/(n(b-a))$ on $[a,b]$ is $n$-CM, again for
$n\ge3$~\cite[Theorem~3.4]{PWW13}.
They also suggested~\cite[Section~6]{PWW12} that their methods
might yield the complete mixability of unimodal asymmetric distributions.
Puccetti and Wang~\cite{PW15} gave numerical evidence that the mean
condition is not sufficient for unimodal densities, with truncated
lognormal distributions and $n=3$; an explicit example is in
Theorem~\ref{thm:sharp} below. Open Problem~10 of~\cite{Wan15} asks:
\begin{quote}
Under what extra conditions a unimodal distribution on a bounded interval
is $n$-CM?
\end{quote}
The survey adds, following Bernard, Jiang and Wang~\cite{BJW14}, that the
question is of particular interest when one of the inequalities
in~\eqref{eq:meancond} is an equality.

The two theorems below settle a special case: $n=3$, the boundary mean
$\mu=a+(b-a)/3$, and densities that are polynomials of low degree. By the
affine invariance of complete mixability~\cite[Proposition~2.1(1)]{WW11} we
may take $[a,b]=[0,1]$, so that $\mu=1/3$ and the constant sum is $1$. We
call a density $f$ on $[0,1]$ \emph{unimodal} if there is $m\in[0,1]$ such
that $f$ is nondecreasing on $[0,m]$ and nonincreasing on $[m,1]$; the mode
$m$ may be an endpoint.

\begin{theorem}\label{thm:main}
Let $f$ be a probability density on $[0,1]$ that is a polynomial of degree
at most three, is unimodal, and has mean $1/3$. Then $f$ is $3$-CM.
\end{theorem}

\begin{theorem}\label{thm:sharp}
There is a probability density $p$ on $[0,1]$ that is a polynomial of degree
four, is strictly positive on $[0,1]$, has exactly one local maximum in
$[0,1]$, which lies in its interior (so $p$ is unimodal), and has mean
$1/3$, but is not $3$-CM.
\end{theorem}

The density in Theorem~\ref{thm:sharp} is given in~\eqref{eq:p}. By
reflection $x\mapsto1-x$, both theorems hold at the other boundary mean
$2/3$, and by affine maps at the boundary means of any interval $[a,b]$.

We call a density on $[0,1]$ that is a polynomial of degree at most three a
\emph{cubic density}. None of the sufficient conditions above applies
directly to a unimodal cubic density $g$ with mean $1/3$ and a mode in
$(0,1)$. Such a $g$ is not monotone. It is not symmetric and not bounded
below by $1$, because a density on $[0,1]$ with either property has mean
$1/2$, and it is not concave, because the only concave density on $[0,1]$
with mean $1/3$ is $2(1-x)$ (Remark~\ref{rem:concave}). Two results of Wang
and Wang reduce Theorem~\ref{thm:main} to a single density with an interior
mode: the theorem on monotone densities, and the convexity of the set of
$n$-CM distributions with a given mean~\cite[Proposition~2.1(3)]{WW11}. If
$g$ is a cubic density with mean $1/3$ and $u=g'(0)$, $v=-g'(1)$, $w=g(1)$,
then (Proposition~\ref{prop:decomp})
\[
g=\frac{3v}4A+6wB+\frac u4H,\qquad\frac{3v}4+6w+\frac u4=1,
\]
where
\[
A(x)=\tfrac{16}9-4x^2+\tfrac{20}9x^3,\qquad
B(x)=\tfrac{11}6-5x^2+\tfrac{10}3x^3,\qquad
H(x)=\tfrac43(1-x)^2(1+5x).
\]
The densities $A$ and $B$ are decreasing with mean $1/3$, and $w\ge0$, so
when $u$ and $v$ are nonnegative, $g$ is a convex combination of two
decreasing densities and $H$. The density $H$ increases on $[0,1/5]$ and
decreases on $[1/5,1]$. It is $3$-CM: put the coordinates of a
Dirichlet$(2,2,1)$ vector in uniformly random order, and each coordinate of
the result has density $H$ (Lemma~\ref{lem:H}). The density $H$ is also a
mixture of a concave density on an interval $[0,\beta]$ with $\beta<1$ and a
nonincreasing density on $[0,1]$, both with mean $1/3$
(Proposition~\ref{prop:mixture}). So its complete mixability, and with it
Theorem~\ref{thm:main}, also follows from the identity above, the theorem of
Puccetti, Wang and Wang on concave densities~\cite[Theorem~4.2]{PWW12}, the
theorem of Wang and Wang on monotone densities and convexity; Lemma~\ref{lem:H} gives the
coupling explicitly.

Every $3$-CM density $f$ on $[0,1]$ with mean $1/3$ that is continuous at
$0$ and $1$ satisfies $f(0)\ge2f(1)$. This follows from the norm condition
of Wang~\cite[Theorem~5]{Wan15} (see~\eqref{eq:wang}), so Wang's theorem
already shows that a unimodal density of mean $1/3$, continuous at $0$ and
$1$, with $f(0)<2f(1)$ is not $3$-CM. Theorem~\ref{thm:sharp} supplies such
a density that is a polynomial of degree four and strictly positive on
$[0,1]$. Lemma~\ref{lem:necessary} proves $f(0)\ge2f(1)$ directly, together
with $f(0)\ge4f(1)$, which improves it by a factor of two. We call these two
inequalities the \emph{endpoint conditions}. The quadratic density
$r(x)=3-8x+6x^2$ satisfies the first and violates the second. So the unimodality assumption in Theorem~\ref{thm:main} cannot
be dropped; see Section~\ref{sec:examples}.

\section{Known results}

\begin{lemma}[{\cite[Theorem~2.4]{WW11}}]\label{lem:monotone}
If $f$ is a nonincreasing probability density on $[0,1]$ with mean at least
$1/n$, then $f$ is $n$-CM.
\end{lemma}

In~\cite{WW11} the theorem is stated for a decreasing density, in the weak
sense: for $n=2$ it covers the uniform density.

\begin{lemma}[{\cite[Proposition~2.1(3)]{WW11}}]\label{lem:convex}
If $F$ and $G$ are $n$-CM, both with mean $\mu$, then $\lambda
F+(1-\lambda)G$ is $n$-CM with mean $\mu$ for every $\lambda\in[0,1]$.
\end{lemma}

In~\cite{WW11} the lemma is stated for distributions with the same center,
the constant sum divided by $n$; for an $n$-CM distribution with finite mean
the center is the mean~\cite[Proposition~2.1(2)]{WW11}. One proof of
Lemma~\ref{lem:convex} is to choose between the two couplings by an
independent coin; the sum is $n\mu$ in both cases. By induction the lemma
extends to finite convex combinations.

\begin{remark}\label{rem:concave}
A concave density $f$ on $[0,1]$ with mean $1/3$ equals $2(1-x)$. For
$0\le x<y<1$ concavity gives
$\frac{f(1)-f(x)}{1-x}\ge\frac{f(1)-f(y)}{1-y}$, so
\[
\varphi(x)=\frac{f(x)}{1-x}=\frac{f(1)}{1-x}-\frac{f(1)-f(x)}{1-x}
\]
is nondecreasing on $[0,1)$, because $f(1)\ge0$. The function
$\omega(x)=(x-\frac13)(1-x)$ has $\int_0^1\omega=0$, is negative on
$(0,1/3)$ and positive on $(1/3,1)$. Hence
\[
0=\int_0^1\Bigl(x-\frac13\Bigr)f(x)\,dx
=\int_0^1\omega(x)\bigl(\varphi(x)-\varphi(\tfrac13)\bigr)\,dx,
\]
and the integrand is nonnegative. So $\varphi$ is constant on $(0,1)$, and
$f=c(1-x)$ with $c=2$ because $\int_0^1f=1$.
\end{remark}

\section{Proof of Theorem~\ref{thm:main}}\label{sec:proof}

Every polynomial $g$ of degree at most three with $\int_0^1g=1$ and
$\int_0^1xg(x)\,dx=1/3$ can be written, with two free real parameters $c$
and $d$, as
\begin{equation}\label{eq:general}
g(x)=2+\frac c6+\frac d5-\Bigl(2+c+\frac{9d}{10}\Bigr)x+cx^2+dx^3.
\end{equation}
Set
\begin{equation}\label{eq:uvw}
u=g'(0)=-2-c-\frac{9d}{10},\qquad v=-g'(1)=2-c-\frac{21d}{10},\qquad
w=g(1)=\frac c6+\frac{3d}{10}.
\end{equation}
Consider the polynomials $A$, $B$ and $H$ of the introduction. Each of $A$, $B$ and $H$ has
integral $1$ and mean $1/3$. Their derivatives are
\[
A'(x)=-\tfrac43x(6-5x),\qquad B'(x)=-10x(1-x),\qquad
H'(x)=4(1-x)(1-5x).
\]
So $A$ and $B$ are decreasing on $[0,1]$, with $A(1)=0$ and $B(1)=1/6$;
in particular they are nonnegative. The polynomial $H$ is nonnegative on
$[0,1]$, increases on $[0,1/5]$ and decreases on $[1/5,1]$.

\begin{lemma}\label{lem:H}
The density $H$ is $3$-CM.
\end{lemma}

\begin{proof}
Let $T=\{(x,y):x,y\ge0,\ x+y\le1\}$, write $z=1-x-y$, and let $(X,Y)$ have
density
\[
q(x,y)=8(xy+xz+yz)
\]
on $T$. This function is nonnegative on $T$ and has integral $1$. It is a
symmetric function of the three barycentric coordinates $x,y,z$. Hence
$X$, $Y$ and $Z=1-X-Y$ have the same distribution. For $0\le x\le1$,
\[
\int_0^{1-x}q(x,y)\,dy=8x(1-x)^2+\frac43(1-x)^3=H(x),
\]
so this common distribution has density $H$. Since $X+Y+Z=1$, $H$ is $3$-CM.
\end{proof}

The density $q$ is $\tfrac13$ times the sum of the three Dirichlet$(2,2,1)$
densities obtained by permuting the parameters, so $(X,Y,Z)$ is a
Dirichlet$(2,2,1)$ vector with its coordinates put in uniformly random
order.

\begin{proposition}\label{prop:mixture}
Let $\beta$ be the root in $(9/10,19/20)$ of
\[
Q(\sigma)=4375\sigma^4-11000\sigma^3+9425\sigma^2-3240\sigma+432;
\]
it is the only root there, and $\beta\approx0.90702$. Then $H=\lambda K+(1-\lambda)D$ for
some $\lambda\in(0,1)$, where $K$ is a concave probability density on $[0,\beta]$ and
$D$ is a nonincreasing probability density on $[0,1]$, both with mean $1/3$.
Consequently $H$ is $3$-CM by~\cite[Theorem~4.2]{PWW12},
Lemma~\ref{lem:monotone} and Lemma~\ref{lem:convex}.
\end{proposition}

\begin{proof}
We have $H'(x)=4(1-x)(1-5x)$ and $H''(x)=8(5x-3)$. For
$\sigma\in[9/10,19/20]$ put $x_\sigma=6/5-\sigma\in[1/4,3/10]$. Then
$H'(x)-H'(\sigma)=4(x-\sigma)(5x+5\sigma-6)$, so $H'(x_\sigma)=H'(\sigma)<0$ and $H'(x)\le H'(\sigma)$ for
$x\in[x_\sigma,\sigma]$. Define
\[
\kappa_\sigma(x)=
\begin{cases}
H(x)-H(x_\sigma)-H'(\sigma)(\sigma-x_\sigma), & 0\le x\le x_\sigma,\\
-H'(\sigma)(\sigma-x), & x_\sigma\le x\le \sigma,\\
0, & \sigma<x\le1,
\end{cases}
\]
and $\delta_\sigma=H-\kappa_\sigma$. The two formulas for $\kappa_\sigma$ agree at $x_\sigma$,
and so do their derivatives, both equal to $H'(\sigma)$. Since $H''<0$ on
$[0,3/5)\supset[0,x_\sigma]$, the function $\kappa_\sigma$ is concave on $[0,\sigma]$. It
vanishes at $\sigma$, and $\kappa_\sigma(0)>0$: as a polynomial in $\sigma$, $\kappa_\sigma(0)$
equals $27/50$ at $\sigma=9/10$ and has no root in $[9/10,19/20]$. So
$\kappa_\sigma>0$ on $[0,\sigma)$. The function $\delta_\sigma$ is constant on $[0,x_\sigma]$,
has derivative $H'(x)-H'(\sigma)\le0$ on $[x_\sigma,\sigma]$, and equals $H$ on $[\sigma,1]$,
where $H$ is nonincreasing. Being continuous, $\delta_\sigma$ is nonincreasing on
$[0,1]$, and $\delta_\sigma\ge H(1)=0$.

Exact integration gives
\[
\int_0^1x\kappa_\sigma(x)\,dx-\frac13\int_0^1\kappa_\sigma(x)\,dx
=-\frac{2(5\sigma-1)Q(\sigma)}{9375}.
\]
Since $Q(9/10)=27/16$ and $Q(19/20)=-1945/256$, and $Q$ has exactly one root
in $[9/10,19/20]$, the root $\beta$ exists and is unique there, and $\kappa_\beta$
has mean $1/3$ after normalization. The integral $\lambda=\int_0^1\kappa_\beta$ is a
value of a polynomial in $\sigma$ that equals $999/2000$ at $\sigma=9/10$ and
$1891/6400$ at $\sigma=19/20$ and takes neither the value $0$ nor the value $1$
on $[9/10,19/20]$; so $0<\lambda<1$. Put $K=\kappa_\beta/\lambda$ and $D=\delta_\beta/(1-\lambda)$.
Then $K$ is a concave density on $[0,\beta]$ with mean $1/3$. Since $H$ has mean
$1/3$, so does $D$, and $D$ is nonincreasing on $[0,1]$. By~\cite[Theorem~4.2]{PWW12}
$K$ is $3$-CM, by Lemma~\ref{lem:monotone} $D$ is $3$-CM, and by
Lemma~\ref{lem:convex} so is $H=\lambda K+(1-\lambda)D$.
\end{proof}

\begin{proposition}\label{prop:decomp}
For $g$ as in~\eqref{eq:general} and $u,v,w$ as in~\eqref{eq:uvw},
\begin{equation}\label{eq:decomp}
g=\frac{3v}4A+6wB+\frac u4H,\qquad\frac{3v}4+6w+\frac u4=1.
\end{equation}
If $u$, $v$ and $w$ are nonnegative, then $g$ is a probability density and
is $3$-CM.
\end{proposition}

\begin{proof}
Both identities in~\eqref{eq:decomp} follow by comparing coefficients
of~\eqref{eq:general}. If $u,v,w\ge0$, then~\eqref{eq:decomp} writes $g$ as
a convex combination of $A$, $B$ and $H$, which are $3$-CM with mean
$1/3$ by Lemmas~\ref{lem:monotone} and~\ref{lem:H}. Lemma~\ref{lem:convex}
completes the proof.
\end{proof}

\begin{proof}[Proof of Theorem~\ref{thm:main}]
Let $f$ be as in the theorem, with mode $m$. If $f$ is constant, its mean
is $1/2$; so $f$ is not constant, and being a polynomial it is strictly
monotone on each of $[0,m]$ and $[m,1]$ that has positive length.

Suppose $m=1$, so that $f$ is nondecreasing on $[0,1]$. Then $f(x)\ge
f(1-x)$ for $x\ge1/2$, and
\[
\int_0^1\Bigl(x-\frac12\Bigr)f(x)\,dx
=\int_{1/2}^1\Bigl(x-\frac12\Bigr)\bigl(f(x)-f(1-x)\bigr)\,dx\ge0,
\]
so the mean is at least $1/2$, a contradiction.

If $m=0$, then $f$ is decreasing on $[0,1]$ with mean $1/3$, and it is
$3$-CM by Lemma~\ref{lem:monotone}.

If $0<m<1$, take $g=f$ in~\eqref{eq:uvw}. Since $f$ is nondecreasing near
$0$ and nonincreasing near $1$, we have $u=f'(0)\ge0$ and $v=-f'(1)\ge0$.
Also $w=f(1)\ge0$ because $f$ is a density. Proposition~\ref{prop:decomp}
shows that $f$ is $3$-CM.
\end{proof}

\begin{remark}
The densities of Proposition~\ref{prop:decomp} are themselves unimodal.
Indeed, if such a $g$ were not unimodal, there would be points $x_1<x_2$ in
$[0,1]$ with $g'(x_1)<0<g'(x_2)$; otherwise $g$ is nondecreasing up to the
infimum of the points where $g'<0$ and nonincreasing after it. Since
$g'(0)\ge0\ge g'(1)$, the polynomial $g'$ would vanish in $[0,x_1)$, in
$(x_1,x_2)$ and in $(x_2,1]$. It has degree at most two, so $g'=0$ and $g$ would
be constant, with mean $1/2$. Together with the proof above, this
shows that a cubic density with mean $1/3$ is unimodal exactly when it is
decreasing or satisfies $u,v,w\ge0$. The cubic densities with mean $1/3$
that are not unimodal have $u<0$ or $v<0$. The quadratic $r$ of
Section~\ref{sec:examples} is one of them, and it is not $3$-CM.
\end{remark}

\section{Two endpoint conditions}

\begin{lemma}\label{lem:necessary}
Let $X_1,X_2,X_3$ be random variables with values in $[0,1]$, each with
distribution $F$, such that $X_1+X_2+X_3=1$ almost surely. Let $X$
have distribution $F$.
\begin{enumerate}
\item For $0<t<1/2$,
\begin{equation}\label{eq:nec1}
\PP(X<t)\ge2\,\PP(X>1-t).
\end{equation}
\item For $0<t<1/4$,
\begin{equation}\label{eq:nec2}
\E(t-X)_+\ge\E\bigl(X-(1-2t)\bigr)_+.
\end{equation}
\end{enumerate}
If $F$ has a density $f$ that is continuous at $0$ and at $1$, then
$f(0)\ge2f(1)$ and $f(0)\ge4f(1)$.
\end{lemma}

\begin{proof}
(1) Since $t<1/2$, at most one of the $X_i$ exceeds $1-t$. If $X_i>1-t$,
the other two sum to $1-X_i<t$, so both are less than $t$. Hence
\[
\sum_{j=1}^3\mathbf 1\{X_j<t\}\ge2\sum_{i=1}^3\mathbf 1\{X_i>1-t\}
\]
pointwise. Taking expectations gives $3\PP(X<t)\ge6\PP(X>1-t)$.

(2) Since $t<1/4$, at most one of the $X_i$ exceeds $1-2t$. If none does,
the right-hand side of
\[
\sum_{j=1}^3(t-X_j)_+\ge\sum_{i=1}^3\bigl(X_i-(1-2t)\bigr)_+
\]
is $0$. If $X_i>1-2t$ and $j,k$ are the other two indices, then
\[
(t-X_j)_++(t-X_k)_+\ge2t-X_j-X_k=2t-1+X_i=X_i-(1-2t).
\]
So the displayed inequality holds pointwise, and taking expectations
gives~\eqref{eq:nec2}.

If $f$ is continuous at the endpoints, then as $t\downarrow0$,
\[
\PP(X<t)=f(0)t+o(t),\qquad\PP(X>1-t)=f(1)t+o(t),
\]
\[
\E(t-X)_+=\int_0^t(t-x)f(x)\,dx=f(0)\frac{t^2}2+o(t^2),
\]
\[
\E\bigl(X-(1-2t)\bigr)_+=\int_{1-2t}^1\bigl(x-(1-2t)\bigr)f(x)\,dx
=f(1)\frac{(2t)^2}2+o(t^2).
\]
Dividing~\eqref{eq:nec1} by $t$ and~\eqref{eq:nec2} by $t^2/2$ and letting
$t\downarrow0$ gives $f(0)\ge2f(1)$ and $f(0)\ge4f(1)$.
\end{proof}

Lemma~\ref{lem:necessary} holds for any three random variables with values
in $[0,1]$, the same distribution and sum $1$; they need not be
exchangeable. Since $f(0)\ge4f(1)$ implies
$f(0)\ge2f(1)$ when $f(1)\ge0$, the second endpoint condition is the
stronger one.

A law-determined norm is a map from random variables to $[0,\infty]$ that
is a norm, takes the same value on identically distributed variables and
is monotone on nonnegative variables~\cite[Definition~5]{Wan15}.
Wang~\cite[Theorem~5]{Wan15} proved that if $F$ is $n$-CM with mean $\mu$
and $X\sim F$, then $\|(X-s)_+\|\le(n-1)\|(X-s')_-\|$ for every
law-determined norm $\|\cdot\|$, every $s\in\mathbb{R}$ and
$s'=(n\mu-s)/(n-1)$. Take the $L^1$ norm, $n=3$, $\mu=1/3$ and $s=1-2t$, so
that $s'=t$. This gives
\begin{equation}\label{eq:wang}
\E\bigl(X-(1-2t)\bigr)_+\le2\,\E(t-X)_+,
\end{equation}
and, as in the proof above, $f(0)\ge2f(1)$. Inequality~\eqref{eq:nec2}
removes the factor $2$ in this case, and the example $r$ of
Section~\ref{sec:examples} satisfies~\eqref{eq:wang} for every
$t\in(0,1/4)$ but not~\eqref{eq:nec2}. In the proof of~\eqref{eq:nec2} the
bound $(X_i-(1-2t))_+\le\sum_{j\ne i}(t-X_j)_+$ is needed for one index $i$
only, since for $t<1/4$ at most one coordinate exceeds $1-2t$; summed over
all three indices it gives only~\eqref{eq:wang}.

\section{Examples}\label{sec:examples}

Both examples below have mean $1/3$. The first is unimodal of degree four
and fails both endpoint conditions; the second is quadratic, not unimodal,
and fails only the second endpoint condition.

\subsection{A quartic that is not \texorpdfstring{$3$}{3}-CM}

Let
\begin{equation}\label{eq:h}
p_0(x)=\frac x{101}\bigl(2360-8070x+9360x^2-3645x^3\bigr).
\end{equation}
Then $\int_0^1p_0=1$, $\int_0^1xp_0(x)\,dx=1/3$, and
\begin{equation}\label{eq:hprime}
p_0'(x)=-\frac{20}{101}(9x-2)P(x),\qquad P(x)=81x^2-138x+59.
\end{equation}
Since $P(x)=81\bigl(x-\tfrac{23}{27}\bigr)^2+\tfrac29>0$, the polynomial $p_0$
increases on $[0,2/9]$ and decreases on $[2/9,1]$. With $p_0(0)=0$ and
$p_0(1)=5/101$, this shows that $p_0$ is nonnegative on $[0,1]$, so it is a
unimodal probability density with mean $1/3$. It violates $f(0)\ge2f(1)$,
so it is not $3$-CM by Lemma~\ref{lem:necessary}, or by~\eqref{eq:wang}.

To obtain a strictly positive example, mix in a small amount of the
decreasing density $2(1-x)$, which has mean $1/3$:
\begin{equation}\label{eq:p}
p(x)=\frac{99}{100}p_0(x)+\frac1{100}\cdot2(1-x).
\end{equation}

\begin{proof}[Proof of Theorem~\ref{thm:sharp}]
The polynomial $p$ has degree four, integral $1$ and mean $1/3$. On
$[2/9,1]$ both $p_0$ and $2(1-x)$ are decreasing, so $p'<0$ there. On
$[0,2/9]$ we have $P>0$ and $P'(x)=162x-138<0$, so
\[
p_0''(x)=-\frac{20}{101}\bigl(9P(x)+(9x-2)P'(x)\bigr)<0,
\]
and $p'$ is strictly decreasing on $[0,2/9]$. As $p'(0)=116719/5050>0$
and $p'(2/9)=-1/50<0$, the polynomial $p$ has exactly one critical point
in $[0,1]$, a maximum in $(0,2/9)$, and it is unimodal. Its minimum over
$[0,1]$ is therefore attained at an endpoint, and
\[
p(0)=\frac1{50},\qquad p(1)=\frac{99}{2020},
\]
so $p>0$ on $[0,1]$. Finally $2p(1)=99/1010>1/50=p(0)$, so $p$ violates
$f(0)\ge2f(1)$ and is not $3$-CM by Lemma~\ref{lem:necessary}, or
by~\cite[Theorem~5]{Wan15} through~\eqref{eq:wang}.
\end{proof}

The violation is visible at a fixed $t$: for $t=1/200$, the difference
$\PP(X<t)-2\PP(X>1-t)$ under $p$ is about $-1.72\times10^{-4}$.

\subsection{Unimodality cannot be dropped}

The polynomial
\[
r(x)=3-8x+6x^2=6\Bigl(x-\frac23\Bigr)^2+\frac13
\]
is a strictly positive density on $[0,1]$ with integral $1$ and mean
$1/3$. It decreases on $[0,2/3]$ and increases on $[2/3,1]$, so it is not
unimodal. It satisfies $r(0)=3\ge2r(1)=2$, but $r(0)=3<4r(1)=4$, so it is
not $3$-CM by Lemma~\ref{lem:necessary}. For $0<t<1/4$ the two sides
of~\eqref{eq:nec2} differ by
\[
\E(t-X)_+-\E\bigl(X-(1-2t)\bigr)_+=-\frac{t^2(3t-1)(5t-1)}2,
\]
which is negative for $0<t<1/5$. For $r$ we have $u=r'(0)=-8<0$, so
Proposition~\ref{prop:decomp} does not apply.

\section{Remarks}\label{sec:remarks}

At the boundary mean $1/n$ with $n\ge4$ the analogue of
Theorem~\ref{thm:main} holds for a simpler reason. For a cubic density $g$
on $[0,1]$ with mean $\mu$, let $u=g'(0)$, $v=-g'(1)$ and $w=g(1)$;
comparing coefficients gives $u/4+3v/4+6w=30\mu-9$. A unimodal cubic
density with mode in $(0,1)$ has $u,v,w\ge0$, so its mean is at least
$3/10>1/4$. Hence a unimodal cubic density with mean $1/n\le1/4$ has mode
$0$ (mode $1$ forces mean at least $1/2$, as in the proof of
Theorem~\ref{thm:main}), and it is $n$-CM by Lemma~\ref{lem:monotone}. We do
not know whether every unimodal cubic density on $[0,1]$ with mean in
$(1/3,2/3)$ is $3$-CM.

By Lemma~\ref{lem:necessary}, a $3$-CM density on $[0,1]$ with mean $1/3$
that is continuous at the endpoints satisfies $f(0)\ge4f(1)$. For a cubic
$g$ as in~\eqref{eq:general}, comparing coefficients
gives
\[
g(0)-4g(1)=\frac u3+\frac{4v}3+7w,
\]
so the conditions $u,v,w\ge0$, which unimodality forces when the mode is
interior, imply the endpoint condition $g(0)\ge4g(1)$ directly.

\section*{Acknowledgments and funding}

The research presented in this paper was supported by the European Research Council (ERC) under the European Union's Horizon Europe research and innovation programme (grant agreement No.~101041711) and by the Simons Foundation as part of the Collaboration on the Mathematical and Scientific Foundations of Deep Learning. It was also supported by Heights Labs, by Convex Nexus Capital and by the Israel Science Foundation (ISF grants 2258/19 and 4101/25).

This work was done by multiple AI agents, Claude (Anthropic) and Codex (OpenAI), connected to Math Brain~\cite{MathBrainRepo} and working together inside the Brain framework~\cite{BrainFramework}. The method is described in~\cite{MathBrainPaper}.

\section*{Declaration of competing interest}

The author holds equity in Heights Labs and in Convex Nexus Capital and has a management role at Convex Nexus Capital.

\section*{Declaration of generative AI and AI-assisted technologies}

This paper is part of a broader research project on proving open problems with AI systems. The work was done by multiple AI agents, Claude (Anthropic) and Codex (OpenAI), connected to Math Brain~\cite{MathBrainRepo} and working together inside the Brain framework~\cite{BrainFramework}; the method is described in~\cite{MathBrainPaper}. OpenAI's Codex (GPT-6.1) proposed the decomposition of Proposition~\ref{prop:decomp}, the coupling of Lemma~\ref{lem:H}, the proofs of Lemma~\ref{lem:necessary} and the examples of Section~\ref{sec:examples}. A second Codex run audited them. A Claude audit run pointed out the mixture of Proposition~\ref{prop:mixture}. Anthropic's Claude (Opus~5.5) checked every identity in exact arithmetic, read the cited results at their sources, and wrote the text. The theorem on monotone densities and the convexity used here are due to Wang and Wang~\cite{WW11}, and the theorem on concave densities to Puccetti, Wang and Wang~\cite{PWW12}. The question is Open Problem~10 of Wang~\cite{Wan15}. The new results are the identity~\eqref{eq:decomp}, which reduces every unimodal cubic density with mean $1/3$ to $A$, $B$ and $H$, the explicit coupling of Lemma~\ref{lem:H}, inequality~\eqref{eq:nec2} and the examples of Section~\ref{sec:examples}. The author directed the work and is responsible for the content of the paper and for decisions concerning its dissemination.

\section*{Data and code availability}

The script \texttt{check.py}, at \url{https://github.com/yspennstate/unimodal-cubic-complete-mixability/tree/b95af7b0d6c716f6185ed64dac2a622bb0bb3b63}, verifies in exact rational arithmetic the identities~\eqref{eq:general} and~\eqref{eq:decomp}, the marginal in Lemma~\ref{lem:H}, the facts used in the proof of Proposition~\ref{prop:mixture} and the values in Sections~\ref{sec:examples} and~\ref{sec:remarks}.

\begin{thebibliography}{99}

\bibitem{BJW14} C.~Bernard, X.~Jiang and R.~Wang, \emph{Risk aggregation with dependence uncertainty}, Insurance Math. Econom. \textbf{54} (2014), 93--108. doi:10.1016/j.insmatheco.2013.11.005.

\bibitem{GR81} N.~Gaffke and L.~R\"uschendorf, \emph{On a class of extremal problems in statistics}, Math. Operationsforsch. Statist. Ser. Optim. \textbf{12} (1981), 123--135. doi:10.1080/02331938108842712.

\bibitem{PW15} G.~Puccetti and R.~Wang, \emph{Detecting complete and joint mixability}, J. Comput. Appl. Math. \textbf{280} (2015), 174--187. doi:10.1016/j.cam.2014.11.050.

\bibitem{PWW12} G.~Puccetti, B.~Wang and R.~Wang, \emph{Advances in complete mixability}, J. Appl. Probab. \textbf{49} (2012), 430--440. doi:10.1239/jap/1339878796.

\bibitem{PWW13} G.~Puccetti, B.~Wang and R.~Wang, \emph{Complete mixability and asymptotic equivalence of worst-possible VaR and ES estimates}, Insurance Math. Econom. \textbf{53} (2013), 821--828. doi:10.1016/j.insmatheco.2013.09.017.

\bibitem{RU02} L.~R\"uschendorf and L.~Uckelmann, \emph{Variance minimization and random variables with constant sum}, in: Distributions with Given Marginals and Statistical Modelling, Kluwer, Dordrecht, 2002, 211--222. doi:10.1007/978-94-017-0061-0\_22.

\bibitem{KS06} M.~Knott and C.~S.~Smith, \emph{Choosing joint distributions so that the variance of the sum is small}, J. Multivariate Anal. \textbf{97} (2006), 1757--1765. doi:10.1016/j.jmva.2006.01.004.

\bibitem{MathBrainRepo} Y.~Shmalo, \emph{Math Brain}, public repository, 2026, \url{https://github.com/yspennstate/math-brain-public/tree/8c9a6bd478c4ec4082347151b7ae16db50c917dd}.

\bibitem{MathBrainPaper} Y.~Shmalo, \emph{Math Brain: shared, auditable memory for multi-agent mathematical research}, October 2026, doi:10.5281/zenodo.24150837, \url{https://github.com/yspennstate/math-brain-public/tree/main/papers/math-brain-2026-10-24}.

\bibitem{BrainFramework} Y.~Shmalo, \emph{Brain Framework}, public software repository, 2026, \url{https://github.com/yspennstate/brain-framework-public/tree/5629e644cdd510173ae97e9aa92642dabd22dc57}.

\bibitem{Wan15} R.~Wang, \emph{Current open questions in complete mixability}, Probab. Surv. \textbf{12} (2015), 13--32. doi:10.1214/14-PS250.

\bibitem{WW11} B.~Wang and R.~Wang, \emph{The complete mixability and convex minimization problems with monotone marginal densities}, J. Multivariate Anal. \textbf{102} (2011), 1344--1360. doi:10.1016/j.jmva.2011.05.002.

\bibitem{YZ17} C.~Yin and D.~Zhu, \emph{Joint mixability of elliptical distributions and related families}, preprint, arXiv:1706.05499v4, 2023.

\end{thebibliography}

\end{document}
